% ===========================================================================
\chapter{Following the traveller}\label{ch:following}

Printed directions are planned once and read once. A navigator is
different: it rides along. About once a second, the GPS gives it a new
position, a \emph{fix}, and it must work out:

\begin{itemize}[nosep]
  \item how far along the route the traveller has got;
  \item what the next instruction is, and how far away;
  \item how far there is still to go, and how long it will take;
  \item and whether the traveller is still on the route at all.
\end{itemize}

\section{Noisy fixes}

A GPS fix is never exactly right. Out in the open, a phone is usually
within about five metres of the truth; between tall buildings, where the
satellites' signals bounce off walls, it can be out by fifty. A fix taken
while driving down the middle of a road might put the car in a garden.

So a navigator never takes a fix at face value. It assumes the traveller
is on the road, most likely on the route it planned, and moves each fix to
the nearest point of the route. This is called \emph{map matching}. The
distance between the fix and that point, the \emph{offset}, says how well
they match: a few metres is noise, but fifty metres, again and again, means
the traveller is somewhere else.

\section{The route as a line}

For following a route, we want it as a line on the map, with a ruler
along it: for any point of the route, how far along it is. We store, for
each node of the route, the distance from the start of the route to that
node, and the time it takes to get there.

\begin{center}
\begin{tikzpicture}[x=1cm, y=1cm]
  \draw[accent, very thick] (0,0) -- (3,0.6) -- (5,0.2) -- (7.5,1.4);
  \foreach \p/\d in {(0,0)/0, (3,0.6)/310, (5,0.2)/520, (7.5,1.4)/800}
    \fill \p circle (0.06) node[below=2pt, font=\scriptsize] {\d\,m};
  \fill[st] (4,0.4) circle (0.07);
  \node[font=\scriptsize, text=st, above] at (4,0.45) {415\,m along};
\end{tikzpicture}
\end{center}

A point $d$ metres along the route is on the segment whose start is at
most $d$ along and whose end is more than $d$ along. Within that segment,
it is a fraction of the way from start to end, and so is its time: the
time at the start of the segment, plus the same fraction of the segment's
time.

\begin{newfile}{navigation.nex}
-- navigation.nex: following a route as the traveller moves. Knowing how
-- far along the route they are, what to tell them next, and when they
-- have left the route.

intern Directions
intern Helpers

-- Where a position lies relative to a line: the nearest point of the
-- line, how far along the line that point is, and how far the position
-- is from it.
class Line_Match
create
  make(p: Point, d, off: Real) do
    point := p
    along := d
    offset := off
  end
feature
  once point: Point
  -- Metres from the start of the line.
  once along: Real
  -- Metres from the line.
  once offset: Real
end

-- A route as a line on the map, measured along its length.
class Route_Line
create
  make(g: Road_Graph, path: Array[Integer])
  require
    is_a_route: is_route(g, path)
  do
    nodes := path
    points := g.island.points_of(path)
    along := [0.0]
    times := [0.0]
    from
      let i: Integer := 1
    until
      i = path.length
    do
      if ?g.edge_between(path.get(i - 1), path.get(i)) as e then
        along.add(along.get(i - 1) + e.length)
        times.add(times.get(i - 1) + e.seconds)
      end
      i := i + 1
    end
  end
feature
  nodes: Array[Integer]
  points: Array[Point]
  -- How far along the route each node is, in metres...
  along: Array[Real]
  -- ...and how long it takes to get there, in seconds.
  times: Array[Real]

  length(): Real do
    result := along.get(along.length - 1)
  end

  duration(): Real do
    result := times.get(times.length - 1)
  end
\end{newfile}

Now the questions about a point $d$ metres along: which segment it is on,
where it is, which way the route is heading there, and how long it takes to
get there. A route with only one node, from a place to itself, has no
segments at all, and each query handles that first.

\begin{continued}{navigation.nex}
  -- The segment that the point `d` metres along the route is on: i for
  -- the segment from node i to node i + 1.
  segment_at(d: Real): Integer do
    result := 0
    from
    until
      result >= points.length - 2 or along.get(result + 1) > d
    do
      result := result + 1
    end
  ensure
    valid_segment: result >= 0 and (result < points.length - 1 or points.length = 1)
  end

  -- The point `d` metres along the route.
  point_at(d: Real): Point do
    if points.length = 1 then
      result := points.get(0)
    else
      let i: Integer := segment_at(d)
      result := points.get(i) + (points.get(i + 1) - points.get(i)) * fraction(i, d)
    end
  end

  -- The direction of travel `d` metres along the route.
  heading_at(d: Real): Real do
    if points.length = 1 then
      result := 0.0
    else
      let i: Integer := segment_at(d)
      result := bearing(points.get(i), points.get(i + 1))
    end
  end

  -- How long it takes to travel the first `d` metres of the route.
  time_at(d: Real): Real do
    if points.length = 1 then
      result := 0.0
    else
      let i: Integer := segment_at(d)
      result := times.get(i) + (times.get(i + 1) - times.get(i)) * fraction(i, d)
    end
  end
\end{continued}

For drawing, the live map will show the part of the route already travelled
in grey and the rest in blue, so the line can be cut in two at any point:

\begin{continued}{navigation.nex}
  -- The route from `d` metres on (or up to `d` metres, if not `after`),
  -- as a line, for drawing.
  part(d: Real, after: Boolean): Array[Point] do
    let i: Integer := segment_at(d)
    if after then
      result.add(point_at(d))
      result := result.concat(points.drop(i + 1))
    else
      result := points.take(i + 1)
      result.add(point_at(d))
    end
  end
\end{continued}

\section{Map matching}

To match a fix to the route, we use \code{closest_point_on_segment}, from
Chapter~\ref{ch:geometry}, to find the closest point of every segment, and
keep the nearest. How far along the route is it? As far as the start of its
segment, plus its distance from that start.

There is a catch. Routes can double back, or run close beside themselves:
out along a lane and back along the same lane, or down one street and up
the next. The nearest point of the whole route might be on a part the
traveller won't reach for minutes, and if we believed it, the navigator
would suddenly announce the instructions for the way back. So we search
only a \emph{window} of the route, between \code{low} and \code{high}
metres along it, around where the traveller was a moment ago. Real
navigators do the same, and more: they also use the direction and speed of
travel, and they weigh up several candidate roads over several fixes before
deciding.

\begin{continued}{navigation.nex}
  -- Where p is relative to the stretch of the route between `low` and
  -- `high` metres along it.
  locate(p: Point, low, high: Real): Line_Match do
    result := create Line_Match.make(points.get(0), 0.0, p.distance_to(points.get(0)))
    from
      let i: Integer := 0
    until
      i >= points.length - 1
    do
      if along.get(i + 1) >= low and along.get(i) <= high then
        let a: Point := points.get(i)
        let c: Point := closest_point_on_segment(p, a, points.get(i + 1))
        let off: Real := p.distance_to(c)
        if off < result.offset then
          result := create Line_Match.make(c, along.get(i) + a.distance_to(c), off)
        end
      end
      i := i + 1
    end
  end

private feature
  -- How far through segment i the point d metres along the route is,
  -- from 0 at its start to 1 at its end.
  fraction(i: Integer, d: Real): Real do
    let seg: Real := along.get(i + 1) - along.get(i)
    if seg > 0.0 then
      result := ((d - along.get(i)) / seg).clamp(0.0, 1.0)
    end
  end
end
\end{continued}

\section{The navigator}

The navigator holds the route and its instructions, and the one number
that matters most: its \emph{progress}, how far along the route the
traveller is. It also keeps the graph and the search algorithm, which it
will need in the next chapter, to plan a new route when the traveller goes
astray.

\begin{append}{navigation.nex}
-- The navigator: it knows the route, takes in GPS fixes, and says what to
-- do next.
class Navigator
create
  make(g: Road_Graph, path: Array[Integer], place: String, how: String)
  require
    is_a_route: is_route(g, path)
  do
    graph := g
    destination := place
    algorithm := how
    progress := 0.0
    misses := 0
    reroutes := 0
    route := create Route_Line.make(g, path)
    instructions := directions(g, path, place)
    last_fix := route.point_at(0.0)
  end
feature
  graph: Road_Graph
  -- The destination's name, and the algorithm used to plan routes to it.
  destination: String
  algorithm: String
  route: Route_Line
  instructions: Array[Instruction]
  -- How far along the route the traveller is, in metres.
  progress: Real
  -- How many fixes in a row have been too far from the route.
  misses: Integer
  -- How many times the route has been planned again.
  reroutes: Integer
  last_fix: Point
\end{append}

Each fix is matched to the route, looking from 40 metres behind the last
progress to 250 metres ahead of it. Driving at 60\,km/h, a car covers about
17 metres a second, so 250 metres is plenty of room ahead.

If the fix matches within 25 metres, the traveller is on the route, and
the progress moves up to where the fix matched. It never goes backwards:
the noise in the fixes makes the matched point jitter back and forth by a
few metres, and a navigator whose ``in 200 m'' becomes ``in 205 m'' looks
broken. The postcondition holds the navigator to that, with \code{old}:
in an \code{ensure}, \code{old progress} is the value \code{progress} had
when the routine was called.

If the fix is further away, it is a \emph{miss}. One miss might be a bad
fix. Three in a row, and the traveller has left the route.

\begin{continued}{navigation.nex}
  -- A fix further than this from the route doesn't count as on it...
  TOLERANCE = 25.0
  -- ...and this many of those in a row mean the traveller has left it.
  PATIENCE = 3
  -- How far back and ahead of the last known progress to look for the
  -- traveller on the route.
  LOOK_BACK = 40.0
  LOOK_AHEAD = 250.0

  -- Take in a new GPS fix.
  --
  -- GPS fixes are a few metres off, so the navigator moves each one to the
  -- nearest point of the route: this is called map matching. It looks
  -- only near where the traveller was a moment ago. A route can pass the
  -- same place twice, or run close to itself, and the nearest point of the
  -- whole route might be on a part the traveller has not reached yet.
  update(fix: Point) do
    last_fix := fix
    let m: Line_Match := route.locate(fix, (progress - LOOK_BACK).max(0.0),
      progress + LOOK_AHEAD)
    if m.offset <= TOLERANCE then
      progress := progress.max(m.along)
      misses := 0
    else
      misses := misses + 1
    end
  ensure
    no_going_back: progress >= old progress
  end

  is_off_route(): Boolean do
    result := misses >= PATIENCE
  end

  has_arrived(): Boolean do
    result := progress >= route.length() - 15.0
  end
\end{continued}

\section{What to say}

The next instruction is the first one further along the route than the
traveller. What to say about it depends on how far away it is: far off,
``In 300 m, turn left onto Mill Lane''; close by, ``Turn left onto Mill
Lane now''. And the time still to go is the route's whole duration, less
the time it takes to get to where the traveller is.

\begin{continued}{navigation.nex}
  -- The next instruction the traveller has not reached yet.
  next_instruction(): Instruction do
    result := instructions.get(instructions.length - 1)
    let found: Boolean := false
    across instructions as step do
      if not found and step.distance > progress + 1.0 then
        result := step
        found := true
      end
    end
  end

  -- What to say now: "In 300 m, turn left onto Mill Lane".
  prompt(): String do
    let next: Instruction := next_instruction()
    let d: Real := next.distance - progress
    if has_arrived() then
      result := "You have arrived at " + destination
    elseif d < 30.0 then
      result := next.text + " now"
    else
      result := "In " + distance_text(d) + ", " + lowercase_first(next.text)
    end
  end

  remaining_distance(): Real do
    result := (route.length() - progress).max(0.0)
  end

  remaining_time(): Real do
    result := (route.duration() - route.time_at(progress)).max(0.0)
  end
end

-- s with its first letter in lower case: "Turn left" -> "turn left".
function lowercase_first(s: String): String do
  if s.length = 0 then
    result := s
  else
    result := s.substring(0, 1).to_lower() + s.substring(1, s.length)
  end
end
\end{continued}

The remaining time comes from the route's own times, edge by edge, so it
knows that the rest of the trip is on a slow track or a fast main road.
Real navigators improve on this all the time, with the traveller's actual
speed and with live traffic.

\section{A pretend drive}

We don't have a traveller yet; the next chapter makes one. For now we can
pretend: take points along the route every 100 metres, push each one a few
metres off to the side, the way a GPS would, and give them to the
navigator. \code{direction(heading + 90.0)} is an arrow pointing square to
the right of the route, so adding six metres of it moves a point six metres
to the right.

\begin{tryit}{try_navigator.nex}
intern Navigation

let island: Island_Map := create Island_Map.load(user_path("kestrel_island.map"))
let g: Road_Graph := create Road_Graph.build(island, Mode.Driving)
if ?island.find_place("harbour") as a and ?island.find_place("lighthouse") as b then
  let found: Search_Result := find_route(g, g.nearest_node(island.point(a.node_id)),
    g.nearest_node(island.point(b.node_id)), "astar", false)
  let nav: Navigator := create Navigator.make(g, found.path, b.name, "astar")
  let line: Route_Line := nav.route
  console.print_line("route: " + distance_text(line.length()) + ", " + duration_text(line.duration()))
  let d: Real := 0.0
  let side: Real := 6.0
  from
  until
    d > line.length() + 100.0
  do
    let fix: Point := line.point_at(d) + direction(line.heading_at(d) + 90.0) * side
    nav.update(fix)
    if (d / 100.0).round % 5 = 0 or nav.has_arrived() then
      console.print_line(nav.progress.format(0).pad_start(" ", 5) + " m  " + nav.prompt().pad_end(" ", 42)
        + distance_text(nav.remaining_distance()) + ", " + duration_text(nav.remaining_time()))
    end
    d := d + 100.0
    side := -side
  end

  -- Another drive, which turns off the route 600 m from the start.
  let lost: Navigator := create Navigator.make(g, found.path, b.name, "astar")
  from
    d := 0.0
  until
    d > 600.0
  do
    lost.update(line.point_at(d))
    d := d + 100.0
  end
  repeat 3 do
    lost.update(line.point_at(600.0) + direction(line.heading_at(600.0) + 90.0) * 60.0)
    console.print_line("a fix 60 m from the route: " + lost.misses.to_string + " in a row, off the route: "
      + lost.is_off_route().to_string)
  end
end
\end{tryit}
\tryoutput{try_navigator}

Every fix was six metres to the left or right of the road, and every one
was matched back onto it. The prompts count down to the next change of
road, and the time left comes down as the distance does.

On the second drive, after 600 metres the fixes are suddenly 60 metres
from the route. One or two such fixes could be bad readings; three in a
row mean the traveller has left the route. What then? That is the next
chapter.
